\documentclass[10pt]{article}
\usepackage[margin=0.85in]{geometry}
\usepackage[T1]{fontenc}
\usepackage{lmodern}
\usepackage{microtype}
\usepackage{booktabs,array,tabularx}
\usepackage{enumitem}
\usepackage{xcolor}
\usepackage[colorlinks=true,linkcolor=black,urlcolor=blue!50!black]{hyperref}
\setlist{itemsep=1pt,topsep=2pt,leftmargin=*}
\setlength{\parskip}{4pt}
\setlength{\parindent}{0pt}
\setlength{\emergencystretch}{2em}
\newcolumntype{L}[1]{>{\raggedright\arraybackslash}p{#1}}
\newcolumntype{Y}{>{\raggedright\arraybackslash}X}
\newcommand{\nctx}{365}
\newcommand{\neone}{2,500}
\newcommand{\exactDSFlash}{98}
\newcommand{\bestmodel}{DeepSeek V4 Flash}
\newcommand{\tiedwithbest}{DeepSeek V4 Pro}
\newcommand{\nsigbelow}{three}
\newcommand{\deltamin}{7}
\newcommand{\deltamax}{17}
\newcommand{\exactmin}{81}
\newcommand{\exactmax}{98}
\newcommand{\recallmin}{0.966}
\newcommand{\recallmax}{0.996}
\newcommand{\otherswrong}{25}
\newcommand{\othersmadeup}{0}
\newcommand{\flashmadeup}{1,032}
\newcommand{\flashworstmade}{991}
\newcommand{\flashworstgold}{100}
\newcommand{\medianmissed}{1}
\newcommand{\ninvalid}{7}
\newcommand{\nrefused}{2}
\newcommand{\failDSPro}{4}
\newcommand{\failDSFlash}{2}
\newcommand{\largeHaiku}{53}
\newcommand{\largeMiMoFlash}{34}
\newcommand{\nenum}{100}
\newcommand{\smallmin}{93}
\newcommand{\orpooledanswer}{0.50}
\newcommand{\orpooledanswerci}{0.36--0.60}
\newcommand{\orpooledtable}{0.39}
\newcommand{\orpooledtableci}{0.22--0.56}
\newcommand{\fmtMiMoFlashJson}{86}
\newcommand{\fmtMiMoFlashNl}{73}
\newcommand{\condsizeone}{13}
\newcommand{\condsizetwo}{6.5}
\newcommand{\condsizethree}{3.5}
\newcommand{\ragrecall}{0.869 for all models}
\newcommand{\nscalar}{35}
\newcommand{\scalarerrors}{11}
\newcommand{\scalarunder}{9}
\newcommand{\rsmallDSPro}{1,050}
\newcommand{\rlargeDSPro}{8,832}
\newcommand{\rsmallHaiku}{173}
\newcommand{\rlargeHaiku}{6,442}
\newcommand{\flashshallowexact}{62}
\newcommand{\flashdeepexact}{72}
\newcommand{\etwonote}{DeepSeek V4 Pro's E2 runs are pending (provider usage limits).}
\newcommand{\agreemin}{87}
\newcommand{\agreemax}{100}
\newcommand{\flipmin}{0}
\newcommand{\flipmax}{13}
\newcommand{\reasonshare}{77--97}
\newcommand{\totalcalls}{3,796}
\newcommand{\failedcalls}{1,163}

\pagestyle{plain}

\begin{document}
{\Large\textbf{Benchmark proposal, v2: set completeness in LLMs and search agents}}\par
\vspace{2pt}
Daryl Alfaro \quad\textbullet\quad \today \quad\textbullet\quad draft for discussion; covers deliverables 1 and 4
and supersedes \texttt{docs/proposal.md}\par
\vspace{-2pt}\rule{\linewidth}{0.4pt}

\textbf{Question.} Can an LLM or agent determine and return the \emph{complete} set of relevant entities when the
information exists but is not already organised as a database?

\textbf{Design in brief.} Every question has one gold set, fixed from a pinned data snapshot before any model is
run. The same question and gold set are posed in three access settings: (A) the candidates as a table in the
prompt, (B) the candidates as natural-language documents in the prompt, and (C) a searchable collection the agent
must explore itself. Questions span three condition types (explicit, relational, semantic) and two real domains
(Magic: The Gathering and scientific publications); the synthetic benchmark\_v0 stays as the controlled baseline.
Answer size and candidate size are controlled, so each lost item can be attributed to access, to reasoning, or to
the size of the set.

\section*{1. Starting point}
benchmark\_v0 (500 synthetic questions, five models) established:
\begin{itemize}
\item \textbf{The gap.} On the \nctx{} questions with every record in the prompt, mean recall is
\recallmin--\recallmax, yet only \exactmin--\exactmax\% of answers are exactly the correct set.
\item \textbf{Degradation beyond independent errors.} Each item's odds of being missed rise 1.50$\times$ per
doubling of answer size [1.24, 1.83] and 3.0$\times$ per doubling of table size [2.25, 4.03] (beta-binomial
regression, \texttt{analysis/item\_level.md}). Claude Haiku 5.5 misses no item in answers of 10 or fewer, yet
leaves out at least one item in about half of its 50- and 100-item answers.
\item \textbf{Limits this proposal addresses:} synthetic records, simple attribute matching, answer size partly
confounded with table size, answers of at most 100 items, and no search setting.
\end{itemize}

\section*{2. Task taxonomy}
{\small
\begin{tabularx}{\linewidth}{@{}L{2.5cm}YL{4.6cm}@{}}
\toprule
Dimension & Levels & Why \\
\midrule
Condition type & \textbf{Explicit:} an attribute or text of one entity. \textbf{Relational:} relationships
between entities (joins, set operations, nesting, conditions defined by aggregates). \textbf{Semantic:}
interpreting descriptions or evidence. & The three levels in Dr.\ Wen's note; v0's SQL-like operations become the
relational category. \\
Set determinacy & The size is implied by the question (``the top 10 for each year'') or the set is open. &
60\% of WideSearch's real tasks are open-set, including every task above 240 rows
(\texttt{analysis/widesearch.md}). Manipulated directly: tell the model the count, or not. \\
Answer size $g$ & 5 to 500 items & The main scaling factor; 10\% of WideSearch's open-set tasks need 200+ rows. \\
Candidate size $N$ & Pool size in the prompt (A, B); collection size (C) & Separates set size from haystack size. \\
Output form & Entity set (primary); set of pairs; an aggregate that needs the complete set (a count) & \\
Domain & Synthetic (control); Magic: The Gathering; scientific publications & \\
\bottomrule
\end{tabularx}}

\textbf{Example tasks with real gold-set sizes} (counted on 2026-11-02 through the Scryfall and OpenAlex APIs; these
are candidates for deliverable 2, not yet verified gold sets):

{\small
\begin{tabularx}{\linewidth}{@{}L{2.3cm}L{1.6cm}YL{1.1cm}@{}}
\toprule
Domain & Type & Question & $|$gold$|$ \\
\midrule
MTG & explicit & Every card with the toxic keyword. & 43 \\
MTG & explicit & Every card whose type line includes Phyrexian. & 484 \\
MTG & relational & Every card printed in both New Phyrexia (2011) and Phyrexia: All Will Be One (2023). & 7 \\
MTG & relational & Every card banned in Modern but legal in Legacy. & 39 \\
MTG & relational & Every legendary creature whose colour identity includes white, blue and black. & 131 \\
MTG & semantic & Every card in Phyrexia: All Will Be One that can remove an opponent's permanent. & 49$^*$ \\
MTG & semantic & Every Standard-legal card that destroys or exiles many permanents at once (board wipes). & 109$^*$ \\
Publications & explicit & Every paper by Percy Liang published 2020--2025. & 196 \\
Publications & relational & Every paper co-authored by Percy Liang and Christopher D.\ Manning. & 26 \\
Publications & semantic & Every 2024 paper citing ``Lost in the Middle'' that proposes a new benchmark. & pool to build \\
\bottomrule
\end{tabularx}}\par
{\footnotesize $^*$Scryfall's community function tags (\texttt{otag:removal}, \texttt{otag:board-wipe}): a starting
pool to adjudicate, not a gold set.}

\textbf{Why semantic gold sets need a written definition.} Dr.\ Wen's example, ``all cards related to Phyrexia'',
has no single answer: 484 cards have Phyrexian in the type line, 768 mention Phyrexia or Phyrexian anywhere in their name,
type, rules or flavour text, and 1,286 have art tagged Phyrexian by Scryfall's taggers. Each reading is a valid
task; the vague prompt is not. Every semantic task therefore states its definition, and the model sees it.

\section*{3. Data sources}
{\small
\begin{tabularx}{\linewidth}{@{}L{2.3cm}L{4.1cm}L{4.2cm}Y@{}}
\toprule
Source & Content & Licence or terms (checked 2026-11-02) & Use \\
\midrule
Scryfall bulk data & 33,357 paper cards; rulings; community function tags and art tags; regenerated daily &
Free for research under the Wizards of the Coast Fan Content Policy; may not be ``simply'' republished & Pin one
dated file by hash. Release questions, gold card IDs and a rebuild script, not the data. \\
MTG Wiki & Lore: planes, characters, storylines & CC BY-NC-SA (per Scryfall's terms) & Real documents for lore
tasks in B and C \\
MTGJSON & Card and set data & MIT licence page; its scope over the data files is not explicit & Cross-check \\
OpenAlex & Works, authors, venues, citations, abstracts & CC0. Free API key (\$1/day of usage); quarterly free
snapshot ($\sim$745 GB) & Extract bounded pools through the API; pin by date \\
dblp & Computer-science venues and authors & Not verified (the site blocks automated access) & Cross-check of
venue lists \\
WideSearch & 200 real tasks with gold tables & CC0 & Reference for realistic sizes, not ground truth \\
\bottomrule
\end{tabularx}}

\section*{4. Ground-truth construction}
\begin{enumerate}
\item \textbf{Explicit and relational conditions:} SQL over the pinned snapshot loaded into SQLite, as in v0. Each
question stores its SQL, and the tests re-execute it against the stored answer.
\item \textbf{Semantic conditions: a closed pool, labelled exhaustively.}
  \begin{enumerate}
  \item A high-recall pre-filter (the union of text, tag and set signals) defines a candidate pool.
  \item Two LLM judges label every pool item against the written definition. A person adjudicates every
  disagreement and a random sample of agreements; agreement (Cohen's $\kappa$) is reported per task.
  \item A random audit of items \emph{outside} the pool bounds what the pre-filter missed: no positives among 600
  audited items puts the miss rate below 0.5\% with 95\% confidence.
  \end{enumerate}
\item \textbf{Identity.} Each entity has a canonical ID (Scryfall oracle ID; OpenAlex work ID or DOI). The scorer
maps names to IDs through an alias table (both faces of double-faced cards, title variants) and keeps v0's
near-miss category. Authors in relational tasks are checked by hand: OpenAlex holds a second, one-paper
``Christopher D.\ Manning'' profile, for example.
\item \textbf{Memory.} MTG sets released after a model's training cutoff form a post-cutoff split, on which
recall from memory cannot help.
\end{enumerate}

\section*{5. Information-access settings}
{\small
\begin{tabularx}{\linewidth}{@{}L{2.4cm}YL{3.5cm}L{3.2cm}@{}}
\toprule
Setting & The model receives & Isolates & Main measure \\
\midrule
P: memory (reference) & The question only & Recall from memory & As v0's memory slice; MTG only \\
A: full context & A pool of $N$ candidates as a table (CSV) & Identification and exhaustive output & As v0 \\
B: distributed evidence & The same pool as natural-language documents. Facts a relational condition needs sit on
different pages (card page, set page, lore page). & Reading and combining unstructured evidence & B$-$A gap on
identical questions \\
C: search (controlled) & The full collection (all card, set and lore pages; a publication pool of similar size)
behind \texttt{search(query)}$\to$10 snippets and \texttt{read(id)}, with a budget of 10, 30 or 100 calls. The agent
decides when to stop. & Discovery and stopping & Where items are lost (\S6) \\
C-web (later) & Live web search & Realism & Same questions; gold from the snapshot \\
\bottomrule
\end{tabularx}}

All settings share the gold sets. In A and B the pool is the gold set plus distractors up to $N$ (kept under about
150k tokens); in C it is the whole collection. Pages for B and C are rendered from the snapshot with fixed
templates, so every fact a question depends on is on some page; real MTG Wiki pages are added for lore tasks.
This is the ``controlled, searchable collection with hidden ground truth'' that Dr.\ Wen's note suggests starting from.

\section*{6. Metrics}
\begin{itemize}
\item \textbf{Primary:} exact-set accuracy, with precision, recall, F1 and v0's error types (missing, near-miss,
wrong, made-up, duplicate). Truncated, refused and unparseable answers count as incorrect and are reported
separately.
\item \textbf{Size curves:} exact-set accuracy and per-item miss rate against $g$ at fixed $N$, with the
independent-error baseline $(1-q)^g$ and the odds ratio per doubling.
\item \textbf{Where items are lost (C):} each gold item is never retrieved, retrieved but omitted, or returned;
wrong inclusions are counted separately. \emph{Evidence coverage} (share of gold items whose evidence the agent
read) is compared with \emph{final coverage} (recall). A run \emph{stopped early} if the agent ended with budget
left and recall below 1.
\item \textbf{Effort:} completeness against the search budget (C) and against reasoning effort (A, B).
\item \textbf{Self-assessment:} the model also states how many items it believes exist. The gap to $|$gold$|$ shows
whether it knows its answer is incomplete.
\item \textbf{Comparisons:} by condition type, setting, domain and representation, paired on identical questions
(bootstrap over questions).
\end{itemize}

\section*{7. Answer-size experiment (deliverable 4)}
\begin{itemize}
\item \textbf{Design:} fully crossed $g \in \{5, 10, 20, 50, 100, 200, 500\} \times N \in \{1{,}000, 2{,}000,
4{,}000\}$. All 21 cells are valid (0.1--50\% of rows match), so answer size and table size are no longer
confounded. Synthetic tables from the v0 generator, in CSV.
\item \textbf{Output limits:} answers are row IDs, about 3,000 output tokens for 500 items. The output cap is 64k
tokens, and truncation is reported, never scored as omission. A names variant at $N = 2{,}000$ links back to v0.
\item \textbf{Seeds and repeats:} $k$ independent tables per cell (seeds; $k$ from the power analysis below), with matching rows at random, recorded
positions. A second sample on a third of the cells measures run-to-run variation.
\item \textbf{Analysis:} beta-binomial regression of each answer's missed count on $\log_2 g$, $\log_2 N$ and row
position, which lets misses clump within an answer. The independent-error null is an odds ratio of 1 for $g$. The
observed share of complete answers is also compared with $(1-q)^g$, with $q$ estimated at the same $N$ from
$g \le 10$.
\item \textbf{Power} (\texttt{analysis/power.md}; 300 simulated experiments per setting, calibrated to v0). The
observed effect, an odds ratio of about 1.5 per doubling, is detected in every setting even with 3 tables per
cell. The table gives power for a smaller effect, 1.2 per doubling:\par
{\small\begin{tabular}{@{}llccc@{}}
\toprule
Miss rate at $g = 50$ & Spread across answers & 3 tables per cell & 5 tables (105 answers) & 10 tables \\
\midrule
1\% & low ($\sigma = 0.5$) & 65\% & 89\% & 100\% \\
1\% & high ($\sigma = 1$) & 53\% & 68\% & 94\% \\
2\% & low & 74\% & 97\% & 100\% \\
2\% & high & 58\% & 77\% & 97\% \\
\bottomrule
\end{tabular}}\par
\textbf{Recommendation:} 5 tables per cell (105 answers per model), and 10 for a model whose effect looks small. At
the nominal 5\% level the test wrongly rejects a true null 4--9\% of the time, so the final analysis calibrates its
critical value by parametric bootstrap.
\item \textbf{Cost:} a CSV row costs about 22--26 prompt tokens, so $N = 4{,}000$ is about 100k tokens. With
$k = 5$, one model needs 105 prompts and about 6M input tokens.
\item \textbf{Real-domain check:} MTG offers a natural ladder of candidate sizes for setting C: one set
(170--271 cards), Standard (5,164), Pioneer (15,086), all paper cards (33,357).
\end{itemize}

\section*{8. Pilot and next steps}
\begin{itemize}
\item \textbf{Pilot (deliverable 3):} 2 domains $\times$ 3 condition types $\times$ 2 tasks = 12 tasks with
verified gold sets (deliverable 2), in settings A, B and C (controlled, 30-call budget). Three models: Claude Haiku
5.5 and MiMo-V2.6 Flash, which make measurable errors in v0, and DeepSeek V4 Pro, the strongest, to test whether
its ceiling holds outside synthetic tables. Two samples each: 216 runs. Search runs dominate the cost (an
estimated 100k--300k input tokens each).
\item \textbf{Order:} gold sets for the 12 pilot tasks $\to$ the A/B/C harness and pilot $\to$ the crossed size
experiment $\to$ decide which dimensions to scale.
\item \textbf{Shared work with Shiyu:} one natural split is tasks and ground truth vs.\ the harness and search
agent, to agree at our first call.
\end{itemize}

\textbf{Decisions needed.}
\begin{enumerate}
\item Is live web search (C-web) in scope for the first paper, or a follow-up?
\item Who adjudicates semantic labels? The pilot needs a few hundred items checked.
\item The model list and budget for the pilot.
\item An OpenAlex API key (free, with an account) for building publication pools.
\end{enumerate}

\end{document}
